% !TEX root = ../main.tex

\chapter{Introduzione}

\begin{comment}
Ciò che viene scritto in questo enviroment non sarà mostrato nel pdf.
Questo comando è utile per nascondere un intero paragrafo senza eliminarlo dal file sorgente. 
\end{comment}
Alcune linee guida per scrivere la tesi: 1) Ogni immagine/tabella deve avere le seguenti caratteristiche: a) Una caption; b) Deve essere esplicitamente commentata all'interno del testo (ricordatevi che se vi riferite alla figura/tabella inserendo anche il corrispettivo numero e.g., Figura X mostra che... allora la prima lettera deve essere scritta in maiuscolo; E ricordatevi che latex mette a disposizione un comando per citare correttamente (ref) senza dover mettere a mano il numeretto corrispondente) c) Le immagini devono essere rifatte tramite PPT o similari e salvate in pdf prima di inserirle nel file latex, altrimenti rischiate di inserire immagini sgranate. 2) Ricordatevi sempre di citare le fonti e di rifrasare, altrimenti rischiate di inccorrere nel reato di plagio (che può comportare anche la revoca del titolo accademico) e che citare non sminuisce il vostro lavoro, anzi, lo rafforza. 3)L'Abstract deve necessariamente contenere le seuguenti informazioni a) Un breve riasssunto della problematica; b)Una breve descrizione di ciò che è stato fatto; c) I risultati ottenuti.
\begin{comment}
Il progressivo sviluppo tecnologico ha favorito, negli ultimi anni, l'affermazione di architetture basate su microservizi, in contrapposizione ai tradizionali sistemi monolitici legati a risorse locali. L'avvento del cloud computing ha rappresentato una svolta rilevante nel modo in cui applicazioni e servizi vengono progettati e distribuiti. Come ci suggerisce Google da una semplice ricerca , il cloud è "un insieme di server remoti e programmi a cui si accede tramite Internet"\footnote{\url{http://google.com}}, in netta distinzione rispetto al classico utilizzo di risorse locali. Questa flessibilità, unita alla semplicità di accesso e gestione, ha spinto un numero crescente di aziende e sviluppatori ad adottare architetture a microservizi basate su questo paradigma.

Di fronte a questo scenario in continua evoluzione, è emersa la necessità di proteggere i dati sensibili riducendo il più possibile l'esposizione ai rischi. Gli ambienti cloud presentano infatti un elevato livello di complessità, e i criminali informatici sviluppano costantemente nuove tattiche e tecniche per infiltrarsi in tali sistemi. È proprio in risposta a questa esigenza che la comunità globale OWASP, nata con la missione di "creare una community aperta a livello globale per migliorare la sicurezza dei software tramite l'istruzione, gli strumenti appropriati e la collaborazione"\footnote{\url{https://owasp.org/about/}}, si è specializzata anche nella sicurezza delle infrastrutture cloud, rappresentando oggi un punto di riferimento anche nella gestione di questo genere di ambienti. Questa organizzazione no-profit definisce ogni 3-4 ,senza una cadenza specifica,  una propria Top 10 dei dieci rischi più critici che gli ambienti cloud devono affrontare. 
\end{comment}